\documentclass[11pt]{article}
\usepackage[margin=1in]{geometry}
\usepackage{amsmath,amssymb,amsthm}
\usepackage{booktabs}
\usepackage[colorlinks=true,linkcolor=blue,citecolor=blue,urlcolor=blue]{hyperref}

\newtheorem{theorem}{Theorem}[section]
\newtheorem{proposition}[theorem]{Proposition}
\newtheorem{lemma}[theorem]{Lemma}
\newtheorem{corollary}[theorem]{Corollary}
\theoremstyle{definition}
\newtheorem{definition}[theorem]{Definition}
\newtheorem{remark}[theorem]{Remark}
\newtheorem{hypothesis}{Imported input}
\renewcommand{\thehypothesis}{I\arabic{hypothesis}}

\newcommand{\E}{\mathbb{E}}
\newcommand{\Var}{\operatorname{Var}}
\newcommand{\Cov}{\operatorname{Cov}}
\newcommand{\Tr}{\operatorname{Tr}}
\renewcommand{\Re}{\operatorname{Re}}
\newcommand{\lam}{\lambda_{\max}}
\newcommand{\supp}{\operatorname{supp}}

\title{Ordinary susceptibility and the structure factor}\author{}\date{}
\begin{document}
\maketitle

\begin{center}
\fbox{\parbox{0.92\textwidth}{
\textbf{Status: INCOMPLETE.} The target inequality $\Var_{\mu_\beta}(r_f)\le C M_f$ of \cite{cur} is \emph{not} proved here, and no claim below should be read as a
proof of it. What this note proves, with complete proofs, is: (i) the target is
\emph{exactly} equivalent to a uniform bound on the momentum-space symbol (structure factor) of
the root covariance of the \emph{frustration} observable $h_b=\tfrac{\beta}{3}\sum_{i<j}|F_{b,i}-F_{b,j}|^2$,
so that plane-wave marks suffice; (ii) the symbol is bounded on average over the Brillouin zone and the
target can fail on at most a small fraction of momenta; (iii) indicator marks control all marks up to a
logarithm of the support; (iv) a \emph{mixture obstruction}: a law satisfying every imported averaged moment
input, and which is a convex combination of two laws each satisfying the Ward, contact, innovation and
heat inputs, has superextensive ordinary variance, so any proof must use the exact Wilson score at the
coupling of the source in a way that is not convex in the law; (v) the standing regime and the
absolute lower bound $\beta\ge\beta_0$ used in the missing inequality cover the logarithmic continuum
trajectory, while any polynomial lower bound $\beta\ge n^{\kappa}$ does not. The missing
inequality is isolated in Section~\ref{sec:missing} as a single explicit statement (MI), and the
deduction of the target from (MI) is given with all constants. Section~\ref{sec:fixedvolume} adds: (vi) a conditional
theorem giving the target on the fixed reference torus $T^*$ from the complex-source analyticity clause of the
interface I-1, with the exact required strengthening; and (vii) a proof that the volume-uniform form of
the target is an infrared clustering statement, not implied by uniformly bounded correlation functions and not needed
for Osterwalder--Schrader reconstruction, so that the actual requirement of the existence argument is the fixed-volume form.}}
\end{center}

\section{Setting, definitions and the imported inputs}\label{sec:setting}

\subsection{Lattice, law and sources}

Let $n\ge 3$, $\Lambda_n=(\mathbb Z/n\mathbb Z)^4$, $G=SU(N)$ with $N\ge 2$ fixed. Positively oriented links
are $b=(x,\nu)$, $x\in\Lambda_n$, $\nu\in\{1,2,3,4\}$, with origin $o(b)=x$ and direction $\nu_b=\nu$.
Plaquettes are $p=(x;\mu\nu)$, $\mu<\nu$, with boundary links $\partial p=\{(x,\mu),(x+e_\mu,\nu),(x+e_\nu,\mu),(x,\nu)\}$,
and $U_p=U_{x,\mu}U_{x+e_\mu,\nu}U_{x+e_\nu,\mu}^{-1}U_{x,\nu}^{-1}$. Put $s_p=N-\Re\Tr U_p\in[0,2N]$,
$S=\sum_p s_p$, and
\begin{equation}\label{eq:law}
d\mu_\beta=Z_\beta^{-1}e^{-\beta S}\,d\lambda,\qquad \lambda=\text{normalized product Haar measure on }G^{\text{links}}.
\end{equation}
The original invariant product-group metric, orthonormal generators $T_A$ and the Casimir constant $C_F$ are those
fixed by the standing conventions; $X_b^A$ denotes the original left-invariant derivative in link $b$ and color $A$.
The three plane matchings are
$\mathcal A_1=\{12,34\}$, $\mathcal A_2=\{13,24\}$, $\mathcal A_3=\{14,23\}$, and
$S_i=\sum_{p:\ \text{plane}(p)\in\mathcal A_i}s_p$, so $S=S_1+S_2+S_3$. Each matching contains exactly one plane
through each direction, so $F_{b,i}:=(X_b^AS_i)_A$ is the original directional force of the unique plane of
$\mathcal A_i$ containing $\nu_b$; the set $\{F_{b,i}\}_{i=1,2,3}$ coincides with $\{F_{b;\mu}\}_{\mu\neq\nu_b}$ of
\cite[definition of the sources]{cur}. Define, for every positively oriented link $b$,
\begin{equation}\label{eq:rows}
\begin{gathered}
F_b=\sum_{i=1}^3F_{b,i},\qquad D_b=\sum_{i=1}^3|F_{b,i}|^2,\qquad P_b=|F_b|^2,\qquad
\ell_b=\sum_A (X_b^A)^2S,\\
r_b=\beta D_b-\tfrac13\ell_b,\qquad h_b=\beta\big(D_b-\tfrac13 P_b\big),\qquad t_b=h_b-r_b .
\end{gathered}
\end{equation}
Thus $t_b=\tfrac13(\ell_b-\beta P_b)$. For deterministic real marks $f=(f_b)$ put
$r_f=\sum_bf_br_b$, $h_f=\sum_b f_bh_b$, $t_f=\sum_bf_bt_b$, and
\[
M_f=\sum_b f_b^2,\qquad \|f\|_1=\sum_b|f_b|,\qquad b_f=\|f\|_\infty .
\]
These are literally the sources $r_f,h_f,t_f$ of \cite[definition of the sources]{cur} and \cite[definition of the unchanged sources]{bhi}; every original
root, affected neighboring root, pointing shift and Casimir coefficient is kept, because nothing in
\eqref{eq:rows} has been altered. The exact Casimir identity (\cite[Casimir and physical-factor display]{bhi}) reads
\begin{equation}\label{eq:casimir}
\ell_b=C_F\sum_{p\ni b}\Re\Tr U_p=C_F\sum_{p\ni b}(N-s_p),\qquad\text{hence}\qquad |\ell_b|\le 6C_FN=:C_\ell .
\end{equation}

\paragraph{Stencils.} Let $D(b)$ be the set of links lying in some plaquette containing $b$; $|D(b)|=19$ and
$b\in D(b)$. Write $b\sim c$ if $c\in D(b)$; this relation is symmetric (it means that $b,c$ lie in a common
plaquette). Each of $r_b,h_b,t_b$ is a smooth function of $(U_c)_{c\in D(b)}$ only.

\paragraph{Translations.} The law \eqref{eq:law} and the maps $U\mapsto r_{(x,\nu)}(U)$ are covariant under the
translation group $\Lambda_n$: $r_{(x+z,\nu)}(U)=r_{(x,\nu)}(\tau_zU)$, and $\mu_\beta\circ\tau_z^{-1}=\mu_\beta$.
The same holds for $h$ and $t$.

\paragraph{Regime.} Throughout, $(n,\beta)$ ranges over the periodic moment regime of \cite[periodic regime display]{cur}:
\begin{equation}\label{eq:regime}
\mathcal R=\{(n,\beta):\ n\ge3,\ \beta\ge1,\ n\ge c_*\log(2+N\beta)\}.
\end{equation}
Constants $C,c,C_1,\dots$ depend only on $N$, $c_*$, the fixed metric and derivative conventions; never on
$n$, $\beta$, $f$ or $\supp f$. The physical source is $r_H(f)=c_\beta^{-1}a^{-2}r_f$ (\cite[Casimir and physical-factor display]{bhi});
every variance bound below is multiplied by $c_\beta^{-2}a^{-4}$ for it.

\subsection{The imported inputs, with their checked scope}

The following inputs are used as black boxes. Each is stated in exactly the form used below, with its location
and hypotheses.

\begin{hypothesis}[Marked-set exponential moment]\label{I1}
There are $t_*\in(0,1)$ and $K_*<\infty$ such that for all $(n,\beta)\in\mathcal R$ and every set $A$ of
plaquettes,
\[
\E_{\mu_\beta}\exp\Big(\frac{t_*\beta}{6}\sum_{p\in A}s_p\Big)\le K_*^{|A|/6}.
\]
Source: \cite[subsection ``The actual multiplaquette input'']{inn}, which restates the positive transfer
argument of \cite{sbr}, including the $SU(N)$ remark there. The regime \eqref{eq:regime} is exactly where $K_*$ is
uniformly bounded.
\end{hypothesis}

\begin{hypothesis}[Finite-stencil force inequality]\label{I2}
There is $C_X<\infty$ with $\sum_A|X_b^As_p|^2\le C_Xs_p$ pointwise, for every link $b$ and plaquette $p\ni b$.
Source: \cite[after the moment display of the section]{cub}; \cite[proof of the directional-current Bernstein estimate]{dmn}.
\end{hypothesis}

\begin{hypothesis}[Ward concentration]\label{I3}
There is $C_3$ such that for all $n\ge3$, $\beta\ge1$ and all real marks $f$, $\Var_{\mu_\beta}(t_f)\le C_3M_f$.
Source: \cite[theorem on global concentration of the unchanged Ward difference]{wwa}; no regime restriction is needed.
\end{hypothesis}

\begin{hypothesis}[First-layer innovation, variance form]\label{I5}
There is $C_5$ such that for all $(n,\beta)\in\mathcal R$, every link set $\mathcal C$ containing at most one link of
each plaquette, $\eta=U_{\mathcal C^c}$, and every real mark $f$:
$\E_{\mu_\beta}\Var_{\mu_\beta}(r_f\mid\eta)\le C_5M_f$, and the same for $h_f$ and for the Casimir row
$\ell[f]=\sum_bf_b\ell_b$. Source: \cite[corollary on ordinary and physical innovations]{inn}; the Casimir row satisfies the
joint factorial moment hypothesis of \cite{inn} trivially since it is bounded.
\end{hypothesis}

\begin{hypothesis}[Heat-pressure curvature]\label{I6}
For nonnegative marks $f$, with $\psi_f(t)=\log Z_f(t)-\log Z_\beta$ the heat pressure of \cite[subsection ``The unchanged heat-pressure formulation'']{cur},
one has $\psi_f''(0)=\Var_{\mu_\beta}(r_f)+\E_{\mu_\beta}J_f$ and $\E|J_f|\le C_6M_f$ on $\mathcal R$.
Source: \cite[proposition on extensive actual-law acceleration and the heat-curvature identities]{whe}; \cite[heat-pressure display]{cur}.
\end{hypothesis}

\begin{hypothesis}[Equivalence of objectives]\label{I8}
\cite[proposition in subsection ``Equivalence of the remaining quantitative objectives'']{rsc} together with the
uncontrolled-mean display of \cite{cub}, the discarded-norm and paid-comparison displays of \cite{rsc} and the
resolvent displays of \cite{cur}: in the common
regime of the periodic moment input and the embedded-cube hypotheses, uniform extensive control of any one of
(1) $\Var_\mu r_f$, (2) the cube-marginal mean variance, (3) the excess-mean variance (2.42), (4) the resolvent bound
(1.6) for the first coordinate layer, (5) the heat curvature $\psi_f''(0)\le CM_f$ for $f\ge0$, implies control of all.
\end{hypothesis}

The tilt-averaged bound (T17) of \cite[Part~VI]{csj} and the Gaussian reference kernel of
\cite[Part~III, (2.7)--(2.10)]{csj} are cited in Section~\ref{sec:missing} for comparison only.

\subsection{Elementary consequences of \ref{I1}--\ref{I2}}

\begin{lemma}\label{lem:moments}
On $\mathcal R$, with $\sigma_b:=\sum_{p\ni b}s_p$ (six plaquettes):
\begin{enumerate}
\item[(a)] $\E_{\mu_\beta}[\beta s_p]\le C_1:=t_*^{-1}\log K_*$ for every plaquette $p$;
\item[(b)] $D_b\le 2C_X\sigma_b$ and $P_b\le 3D_b\le 6C_X\sigma_b$; hence, with $C_0=\max\{2C_X,2C_FN\}$,
\[
|r_b|\le C_0(1+\beta\sigma_b),\qquad |h_b|\le C_0\beta\sigma_b,\qquad |t_b|\le C_0(1+\beta\sigma_b);
\]
\item[(c)] for every $q\ge1$, $\|r_b\|_{L^q(\mu_\beta)}+\|h_b\|_{L^q(\mu_\beta)}+\|t_b\|_{L^q(\mu_\beta)}\le C'\,q$;
in particular there is $C_2$ with $\Var(r_b)\vee\Var(h_b)\vee\Var(t_b)\le C_2$ and
$|\Cov(r_b,r_c)|\vee|\Cov(h_b,h_c)|\le C_2$ for all $b,c$;
\item[(d)] $0\le \E_{\mu_\beta}h_b\le C_h:=12C_XC_1$.
\end{enumerate}
\end{lemma}
\begin{proof}
(a) Apply \ref{I1} with $A=\{p\}$ and Jensen: $\exp(\tfrac{t_*\beta}{6}\E s_p)\le K_*^{1/6}$.
(b) $F_{b,i}$ is the sum of $X_bs_p$ over the two plaquettes $p\ni b$ of plane $\mathcal A_i$ through $b$, so
$|F_{b,i}|^2\le 2\sum_{p}|X_bs_p|^2\le2C_X\sum_ps_p$ by \ref{I2}; summing over $i$ gives $D_b\le2C_X\sigma_b$.
$P_b=|\sum_iF_{b,i}|^2\le3D_b$ by Cauchy--Schwarz. Then use \eqref{eq:casimir}.
(c) By \ref{I1} with $A=\{p\ni b\}$, $\E e^{t_*\beta\sigma_b/6}\le K_*$, so $\beta\sigma_b$ has all moments with
$\|\beta\sigma_b\|_q\le (6/t_*)qK_*^{1/q}$; combine with (b).
(d) $h_b\ge0$ is Lemma~\ref{lem:frustration} below; $\E h_b\le\beta\E D_b\le 2C_X\E[\beta\sigma_b]\le 12C_XC_1$ by (a).
\end{proof}

\section{The frustration form of the ordinary source and exact locality}\label{sec:frustration}

\begin{lemma}[Frustration identity]\label{lem:frustration}
For every link $b$,
\begin{equation}\label{eq:frustration}
3D_b-P_b=\sum_{1\le i<j\le3}|F_{b,i}-F_{b,j}|^2,\qquad\text{hence}\qquad
h_b=\frac{\beta}{3}\sum_{i<j}|F_{b,i}-F_{b,j}|^2\ \ge 0 .
\end{equation}
Consequently, for $f\ge0$: $h_f\ge0$, $\E_{\mu_\beta}r_f=\E_{\mu_\beta}h_f\in[0,C_h\|f\|_1]$, and pointwise
$r_f\ge-\tfrac13\ell[f]\ge-2C_FN\|f\|_1$.
\end{lemma}
\begin{proof}
$\sum_{i<j}|F_i-F_j|^2=2\sum_i|F_i|^2-2\sum_{i<j}F_i\cdot F_j=2\sum_i|F_i|^2-\big(|\sum_iF_i|^2-\sum_i|F_i|^2\big)
=3\sum_i|F_i|^2-|\sum_iF_i|^2$. The mean statement uses $\E t_f=0$ (\ref{I3}, or the divergence identity
\eqref{eq:div} below) and Lemma~\ref{lem:moments}(d). The pointwise bound uses $D_b\ge0$ and \eqref{eq:casimir}.
\end{proof}

\begin{remark}
Identity \eqref{eq:frustration} says that, up to the Ward divergence $t_f$ (which has extensive variance by
\ref{I3}), the ordinary source is the total \emph{frustration} between the three matching forces at each root.
In the continuum normalization $F_{b,i}\leftrightarrow a^3D_{\nu_i}F_{\nu_i\mu}$ (no sum), so $h_b$ is a
dimension-six density; this is why the Gaussian reference covariance of $r$ decays like $\mathrm{dist}^{-12}$
(\cite[Part~III, (2.7)]{csj}), and it is the quantity whose nonlinear structure factor is at issue below.
\end{remark}

\begin{lemma}[Divergence form of $t_b$ and exact locality of the contacts]\label{lem:locality}
For every link $b$ and every smooth $G$,
\begin{equation}\label{eq:div}
t_b\,e^{-\beta S}=\tfrac13\,\mathrm{div}_\lambda\big(e^{-\beta S}\textstyle\sum_AF_b^AX_b^A\big),\qquad
\Cov_{\mu_\beta}(t_b,G)=-\tfrac13\,\E_{\mu_\beta}\Big[\sum_AF_b^AX_b^AG\Big].
\end{equation}
Hence $\Cov(t_b,t_c)=\Cov(t_b,h_c)=\Cov(t_b,r_c)=0$ whenever $b\not\sim c$, and for all such pairs
\begin{equation}\label{eq:rh}
\Cov_{\mu_\beta}(r_b,r_c)=\Cov_{\mu_\beta}(h_b,h_c)\qquad (b\not\sim c).
\end{equation}
\end{lemma}
\begin{proof}
The left-invariant fields $X_b^A$ are divergence-free for Haar measure, so
$\sum_AX_b^A(e^{-\beta S}F_b^A)=e^{-\beta S}\sum_A(-\beta F_b^AF_b^A+X_b^AF_b^A)=e^{-\beta S}(\ell_b-\beta P_b)=3t_be^{-\beta S}$,
which is the first identity. Integrating $G$ against it and integrating by parts on the compact group gives
$\E[t_bG]=-\tfrac13\E[F_b\cdot X_bG]$; taking $G\equiv1$ shows $\E t_b=0$, which gives the covariance form.
If $b\not\sim c$ then $t_c,h_c,r_c$ do not depend on $U_b$, so $X_bG=0$ for these $G$. Finally
$\Cov(r_b,r_c)=\Cov(h_b-t_b,h_c-t_c)$ and the three terms containing $t$ vanish.
This is the content of \cite[subsection ``Local Ward moments and the finite-range contact matrix'']{wwa}, reproduced for completeness.
\end{proof}

\section{The structure-factor form of the target}\label{sec:symbol}

\begin{definition}\label{def:symbol}
For $\nu,\nu'\in\{1,\dots,4\}$ and $z\in\Lambda_n$ put
\[
c_{\nu\nu'}(z)=\Cov_{\mu_\beta}\big(r_{(x+z,\nu)},r_{(x,\nu')}\big),\qquad
c^h_{\nu\nu'}(z)=\Cov_{\mu_\beta}\big(h_{(x+z,\nu)},h_{(x,\nu')}\big),
\]
which are independent of $x$ by translation covariance. For $p\in\Lambda_n^*=\frac{2\pi}{n}\Lambda_n$ define the
$4\times4$ matrices (\emph{symbols}, or structure factors)
\[
\hat c(p)=\sum_{z\in\Lambda_n}c(z)e^{-ip\cdot z},\qquad \hat c^{\,h}(p)=\sum_zc^h(z)e^{-ip\cdot z}.
\]
For complex marks $f$ set $\Var(r_f):=\E|r_f-\E r_f|^2$. Define the finite-volume susceptibility constants
\[
\Sigma(n,\beta)=\sup_{f\ne0}\frac{\Var_{\mu_\beta}(r_f)}{M_f},\qquad
\Sigma^h(n,\beta)=\sup_{f\ne0}\frac{\Var_{\mu_\beta}(h_f)}{M_f},
\]
the suprema over real marks (equivalently complex marks, see the proof below).
\end{definition}

\begin{theorem}[Exact structure-factor form]\label{thm:symbol}
For every $(n,\beta)$ with $n\ge3$, $\beta\ge1$:
\begin{enumerate}
\item[(a)] Each $\hat c(p)$ is Hermitian positive semidefinite, and
$\hat c(p)_{\nu\nu'}=n^{-4}\,\E\big[\overline{(R_{p,\nu}-\E R_{p,\nu})}\,(R_{p,\nu'}-\E R_{p,\nu'})\big]$ with
$R_{p,\nu}=\sum_{x}e^{ip\cdot x}r_{(x,\nu)}$.
\item[(b)] For every complex mark $f$, with $\hat f(p,\nu)=n^{-2}\sum_xe^{-ip\cdot x}f(x,\nu)$,
\[
\Var_{\mu_\beta}(r_f)=\sum_{p\in\Lambda_n^*}\hat f(p)^*\,\hat c(p)\,\hat f(p),\qquad M_f=\sum_p|\hat f(p)|^2 .
\]
\item[(c)] $\displaystyle \Sigma(n,\beta)=\max_{p\in\Lambda_n^*}\lam\big(\hat c(p)\big)
=\max_{p,\ |v|=1}\ n^{-4}\Var_{\mu_\beta}\Big(\sum_{x,\nu}v_\nu e^{ip\cdot x}r_{(x,\nu)}\Big).$
\end{enumerate}
The same statements hold with $r$ replaced by $h$ (or by $t$), with $\hat c^{\,h}$ and $\Sigma^h$.
\end{theorem}
\begin{proof}
(a) Writing $\tilde R_{p,\nu}=R_{p,\nu}-\E R_{p,\nu}$ and using translation covariance,
\[
\E\big[\overline{\tilde R_{p,\nu}}\,\tilde R_{p',\nu'}\big]=\sum_{x,y}e^{-ip\cdot x}e^{ip'\cdot y}c_{\nu\nu'}(x-y)
=\sum_ze^{-ip\cdot z}c_{\nu\nu'}(z)\sum_ye^{-i(p-p')\cdot y}=n^4\,\delta_{pp'}\,\hat c(p)_{\nu\nu'} .
\]
At $p=p'$ this exhibits $\hat c(p)$ as a Gram matrix, hence Hermitian positive semidefinite.
(b) Fourier inversion gives $f(x,\nu)=n^{-2}\sum_pe^{ip\cdot x}\hat f(p,\nu)$, hence
$r_f-\E r_f=n^{-2}\sum_{p,\nu}\hat f(p,\nu)\tilde R_{p,\nu}$; expand $\E|r_f-\E r_f|^2$ and use the
display in (a) for all pairs $p,p'$. Plancherel for the unitary normalization $n^{-2}$ gives $M_f=\sum_p|\hat f(p)|^2$.
(c) By (b), $\sup_f\Var(r_f)/M_f$ over complex $f$ equals $\max_p\lam(\hat c(p))$, and the supremum is attained at
$f(x,\nu)=v_\nu e^{ip\cdot x}$ with $v$ a top eigenvector; for this $f$, $M_f=n^4|v|^2$, giving the last expression. The quadratic form
$f\mapsto\sum_{b,c}\bar f_bf_cC_{bc}$ with $C_{bc}=\Cov(r_b,r_c)$ real symmetric positive semidefinite has the same
supremum over real and complex vectors (split $f$ into real and imaginary parts, which contribute additively).
\end{proof}

\begin{proposition}[Ward reduction of the symbol]\label{prop:ward-symbol}
For all $(n,\beta)\in\mathcal R$ and all $p\in\Lambda_n^*$,
\[
\hat c(p)=\hat c^{\,h}(p)+\hat e(p),\qquad \|\hat e(p)\|_{\mathrm{op}}\le C_W:=38\,C_2,
\]
where $\hat e$ is the symbol of $E_{bc}=\Cov(r_b,r_c)-\Cov(h_b,h_c)$, a matrix supported on pairs $b\sim c$.
Consequently
\[
|\Sigma(n,\beta)-\Sigma^h(n,\beta)|\le C_W .
\]
\end{proposition}
\begin{proof}
By \eqref{eq:rh}, $E_{bc}=0$ unless $b\sim c$, and $|E_{bc}|\le 2C_2$ by Lemma~\ref{lem:moments}(c). Each row and
column of $E$ has at most $|D(b)|=19$ nonzero entries, so Schur's test gives $\|E\|_{\mathrm{op}}\le 38C_2$.
Since $E$ is translation covariant, its symbol satisfies $\|\hat e(p)\|\le\|E\|_{\mathrm{op}}$ for every $p$ (the
symbol at $p$ is the compression of $E$ to the plane waves at $p$). The last claim follows from
Theorem~\ref{thm:symbol}(c) applied to $r$ and to $h$.
\end{proof}

\begin{corollary}[Plane-wave marks suffice]\label{cor:planewave}
The target $\Var_{\mu_\beta}(r_f)\le CM_f$ for all real marks $f$ and all $(n,\beta)\in\mathcal R'\subset\mathcal R$ holds
if and only if
\[
\sup_{(n,\beta)\in\mathcal R'}\ \max_{p\in\Lambda_n^*,\ |v|=1}\ n^{-4}\,
\Var_{\mu_\beta}\Big(\sum_{x,\nu}v_\nu e^{ip\cdot x}h_{(x,\nu)}\Big)<\infty ,
\]
and then the two constants differ by at most $C_W$.
\end{corollary}

\begin{corollary}[The symbol is bounded on average over momenta]\label{cor:mean}
For all $(n,\beta)\in\mathcal R$,
\[
\frac1{n^4}\sum_{p\in\Lambda_n^*}\Tr\hat c^{\,h}(p)=\sum_{\nu=1}^4\Var_{\mu_\beta}\big(h_{(0,\nu)}\big)\le 4C_2,
\]
hence, for every $K>0$,
\[
\#\{p\in\Lambda_n^*:\ \lam(\hat c^{\,h}(p))>K\}\le \frac{4C_2}{K}\,n^4,
\]
and the same for $\hat c$. In particular the target can fail, for a given $(n,\beta)$ and a given constant $K$,
on at most a fraction $4C_2/K$ of the Brillouin zone.
\end{corollary}
\begin{proof}
$n^{-4}\sum_p\hat c^{\,h}(p)_{\nu\nu}=c^h_{\nu\nu}(0)$ by Fourier inversion, and
$c^h_{\nu\nu}(0)=\Var(h_{(0,\nu)})\le C_2$. Since $\hat c^{\,h}(p)\ge0$, $\lam\le\Tr$, and Markov's inequality
applies to the nonnegative sequence $p\mapsto\Tr\hat c^{\,h}(p)$.
\end{proof}

\begin{proposition}[Summable kernels are sufficient]\label{prop:schur}
If $\sum_{z\in\Lambda_n}\|c^h(z)\|_{\mathrm{op}}\le K_S$ (operator norm of the $4\times4$ block), then
$\Sigma^h(n,\beta)\le K_S$ and $\Sigma(n,\beta)\le K_S+C_W$. Conversely, $\Sigma^h\le C$ for all $(n,\beta)$ does not
imply any decay of $c^h(z)$ in $z$; the two conditions coincide only in the sense of
Theorem~\ref{thm:symbol}(c).
\end{proposition}
\begin{proof}
$\|\hat c^{\,h}(p)\|\le\sum_z\|c^h(z)\|$ for every $p$; apply Theorem~\ref{thm:symbol}(c) and Proposition~\ref{prop:ward-symbol}.
The converse statement is the trivial remark that an operator bound constrains only the symbol, not the
pointwise kernel; see also \cite[Part~III, subsection ``An actual observation can make conditional-mean covariance \dots'']{csj} for an exact example in which a positive
observation destroys pointwise decay while preserving the operator bound.
\end{proof}

\section{Reductions among classes of marks}\label{sec:marks}

\begin{proposition}[Signed marks]\label{prop:signed}
For every real mark $f=f^+-f^-$: $\Var(r_f)\le2\Var(r_{f^+})+2\Var(r_{f^-})$ and
$M_{f^+}+M_{f^-}=M_f$. Hence any bound $\Var(r_g)\le KM_g$ valid for all nonnegative marks $g$ implies
$\Var(r_f)\le 2KM_f$ for all real marks.
\end{proposition}
\begin{proof}
$\Var(X-Y)\le2\Var X+2\Var Y$, and $f^+,f^-$ have disjoint supports.
\end{proof}

\begin{proposition}[Indicator marks control all marks up to a logarithm]\label{prop:indicator}
Fix $(n,\beta)$ and suppose $\Var_{\mu_\beta}(r_{\mathbf 1_A})\le K|A|$ for every set $A$ of links. Then for every
real mark $f$ supported on $m$ links,
\[
\Var_{\mu_\beta}(r_f)\le 2K\,(1+\log m)\,M_f .
\]
The same implication holds for $h$ in place of $r$.
\end{proposition}
\begin{proof}
Let $C_{bc}=\Cov(r_b,r_c)$; $\langle g,Cg\rangle=\Var(r_g)$ is a positive semidefinite form. By
Proposition~\ref{prop:signed} it suffices to treat $u\ge0$ supported on $m$ links with $M_u=\sum u_b^2$.
Layer-cake gives $u=\int_0^\infty\mathbf 1_{\{u>s\}}\,ds$, and by Cauchy--Schwarz for the form $C$,
\[
\begin{aligned}
\langle u,Cu\rangle&=\int_0^\infty\!\!\int_0^\infty\langle\mathbf 1_{\{u>s\}},C\mathbf 1_{\{u>t\}}\rangle\,ds\,dt\\
&\le\Big(\int_0^\infty\langle\mathbf 1_{\{u>s\}},C\mathbf 1_{\{u>s\}}\rangle^{1/2}ds\Big)^2
\le K\Big(\int_0^\infty|\{u>s\}|^{1/2}ds\Big)^2 .
\end{aligned}
\]
Let $u^*_1\ge\dots\ge u^*_m$ be the decreasing rearrangement, $u^*_{m+1}=0$. Then
$\int_0^\infty|\{u>s\}|^{1/2}ds=\sum_{k=1}^m(u^*_k-u^*_{k+1})k^{1/2}=\sum_{k=1}^mu^*_k\big(k^{1/2}-(k-1)^{1/2}\big)
\le\sum_ku^*_kk^{-1/2}\le M_u^{1/2}\big(\sum_{k\le m}k^{-1}\big)^{1/2}\le M_u^{1/2}(1+\log m)^{1/2}$.
Hence $\langle u,Cu\rangle\le K(1+\log m)M_u$, and Proposition~\ref{prop:signed} gives the factor $2$.
\end{proof}

\begin{remark}
For indicator marks $M_f=\|f\|_1=|A|$ and $h_{\mathbf 1_A}\ge0$ (Lemma~\ref{lem:frustration}). The
tilt-averaged estimate (T17) of \cite[Part~VI]{csj} reads for these marks
$\int_{-\delta/2}^{\delta/2}\Var_{\mu_\eta}(r_{\mathbf 1_A})\,d\eta\le C|A|$ with $d\mu_\eta\propto e^{\eta r_{\mathbf 1_A}}d\mu_\beta$,
i.e.\ the target holds for most tilts $\eta$ in a fixed window; it does not give $\eta=0$, exactly as stated there.
Proposition~\ref{prop:indicator} shows that nothing beyond a logarithm of the support is lost by restricting a
future proof to indicator marks.
\end{remark}

\section{A mixture obstruction: the averaged inputs cannot imply the target}\label{sec:mixture}

The following theorem makes precise which of the imported inputs are \emph{convex in the law} and therefore
cannot, by themselves, imply an extensive susceptibility.

\begin{theorem}[Mixture obstruction]\label{thm:mixture}
Let $C_D:=12C_XC_1$ and $\Lambda_0:=2(C_D+2C_FC_1)/(C_FN)$. Fix $\Lambda\ge\max\{\Lambda_0,1\}$ and
$(n,\beta)$ with $n\ge3$, $\beta\ge1$, $n\ge c_*\log(2+N\Lambda\beta)$ (so that both $(n,\beta)$ and
$(n,\Lambda\beta)$ lie in $\mathcal R$). Let
\[
\tilde\mu=\tfrac12\mu_\beta+\tfrac12\mu_{\Lambda\beta}
\]
and keep the sources $r_f,h_f,t_f$ of \eqref{eq:rows} with the coupling $\beta$. Then:
\begin{enumerate}
\item[(i)] $\tilde\mu$ satisfies \ref{I1} with the same $t_*$ and with $K_*$ replaced by
$\max\{K_*(\beta),K_*(\Lambda\beta)\}$, hence every conclusion of Lemma~\ref{lem:moments}, the joint factorial
moment bound (\cite[lemma on joint factorial moments]{inn}) and the tilt-averaged estimate (T17), all with constants depending only on
the stated data.
\item[(ii)] Each component $\mu_\gamma$, $\gamma\in\{\beta,\Lambda\beta\}$, satisfies \ref{I3}, \ref{I5},
\ref{I6} and Lemma~\ref{lem:locality} with its own coupling $\gamma$ and sources; moreover each component satisfies
the variance-form innovation bound \ref{I5} for the $\beta$-sources $r_f,h_f$.
\item[(iii)] Nevertheless, for every nonnegative mark $f$,
\[
\Var_{\tilde\mu}(r_f)\ \ge\ \tfrac{9}{16}\,C_F^2N^2\,\|f\|_1^2,\qquad
\Var_{\tilde\mu}(t_f)\ \ge\ \tfrac14\,C_F^2\big(1-\tfrac1\Lambda\big)^2\big(N-\tfrac{C_1}{\Lambda}\big)^2\|f\|_1^2 .
\]
In particular, for $f=\mathbf 1_A$, $\Var_{\tilde\mu}(r_f)\ge c|A|^2$ while $M_f=|A|$: the target and the Ward
concentration \ref{I3} both fail for $\tilde\mu$ at the coupling $\beta$.
\end{enumerate}
\end{theorem}

\begin{proof}
(i) For any plaquette set $A$ and the $\mu_{\Lambda\beta}$ component, since $s_p\ge0$ and $\Lambda\ge1$,
\[
\E_{\Lambda\beta}\exp\Big(\tfrac{t_*\beta}{6}\sum_{p\in A}s_p\Big)\le\E_{\Lambda\beta}\exp\Big(\tfrac{t_*\Lambda\beta}{6}\sum_{p\in A}s_p\Big)
\le K_*(\Lambda\beta)^{|A|/6};
\]
average the two components. Lemma~\ref{lem:moments}, the joint factorial moment lemma of \cite{inn} and
(T17) of \cite{csj} are derived from \ref{I1} and the pointwise stencil bounds alone, so their proofs apply
verbatim.

(ii) \ref{I3}, \ref{I6} and Lemma~\ref{lem:locality} are theorems about the single Wilson law $\mu_\gamma$ and its own
sources. For \ref{I5} at coupling $\beta$ under $\mu_{\Lambda\beta}$ write
$r_f=\Lambda^{-1}r^{(\Lambda\beta)}_f-(1-\Lambda^{-1})\tfrac13\ell[f]$ with $r^{(\Lambda\beta)}_f=\Lambda\beta D[f]-\tfrac13\ell[f]$,
and use \ref{I5} for $r^{(\Lambda\beta)}_f$ and for the Casimir row, with $\Var(X+Y\mid\eta)\le2\Var(X\mid\eta)+2\Var(Y\mid\eta)$.
The same decomposition handles $h_f$.

(iii) For a two-point mixture, $\Var_{\tilde\mu}X=\tfrac12\Var_\beta X+\tfrac12\Var_{\Lambda\beta}X+\tfrac14(\E_\beta X-\E_{\Lambda\beta}X)^2
\ge\tfrac14(\E_\beta X-\E_{\Lambda\beta}X)^2$. We bound the two means of $r_b$.

Under $\mu_\beta$: $\E_\beta r_b=\E_\beta h_b\ge0$ by Lemma~\ref{lem:frustration}.

Under $\mu_{\Lambda\beta}$: by Lemma~\ref{lem:moments}(b),(a) applied at coupling $\Lambda\beta$,
$\E_{\Lambda\beta}[\beta D_b]=\Lambda^{-1}\E_{\Lambda\beta}[\Lambda\beta D_b]\le\Lambda^{-1}\cdot2C_X\cdot6C_1=C_D/\Lambda$,
and by \eqref{eq:casimir} and Lemma~\ref{lem:moments}(a),
$\tfrac13\E_{\Lambda\beta}\ell_b=\tfrac{C_F}{3}\sum_{p\ni b}(N-\E_{\Lambda\beta}s_p)\ge2C_F\big(N-\tfrac{C_1}{\Lambda\beta}\big)
\ge2C_FN-\tfrac{2C_FC_1}{\Lambda}$. Hence
\[
\E_{\Lambda\beta}r_b\le\frac{C_D+2C_FC_1}{\Lambda}-2C_FN\le-\tfrac32C_FN\qquad(\Lambda\ge\Lambda_0).
\]
Therefore $\E_\beta r_b-\E_{\Lambda\beta}r_b\ge\tfrac32C_FN$ for every $b$, and for $f\ge0$,
$\E_\beta r_f-\E_{\Lambda\beta}r_f\ge\tfrac32C_FN\|f\|_1$, which gives the first bound.

For $t_f$: $\E_\beta t_f=0$ by Lemma~\ref{lem:locality}, while under $\mu_{\Lambda\beta}$ the same divergence identity at
coupling $\Lambda\beta$ gives $\E_{\Lambda\beta}[\Lambda\beta P_b]=\E_{\Lambda\beta}\ell_b$, so
$\E_{\Lambda\beta}t_b=\tfrac13\E_{\Lambda\beta}[\ell_b-\beta P_b]=\tfrac13(1-\Lambda^{-1})\E_{\Lambda\beta}\ell_b
\ge2C_F(1-\Lambda^{-1})(N-C_1/\Lambda)$, using $\beta\ge1$. Then $\E_{\Lambda\beta}t_f\ge2C_F(1-\Lambda^{-1})(N-C_1/\Lambda)\|f\|_1$ for
$f\ge0$ and the mixture variance bound gives the second inequality.
\end{proof}

\begin{corollary}\label{cor:obstruction}
No inequality of the form $\Var(r_f)\le CM_f$ can be deduced from \ref{I1}, \ref{I2}, the local structure of the
rows \eqref{eq:rows}, the tilt-averaged estimate, and the validity of \ref{I3}, \ref{I5}, \ref{I6} and
Lemma~\ref{lem:locality} for the component laws of a convex combination. Any proof of the target must use the
exact Wilson score $\beta\nabla S$ of the single law $\mu_\beta$, at the coupling $\beta$ of the source
\eqref{eq:rows}, through an input that is not stable under convex combinations of laws; the Ward concentration
\ref{I3} is such an input (it fails for $\tilde\mu$), and so is the resolvent identity of \cite[subsection ``The actual retained susceptibility'']{cur}.
\end{corollary}
\begin{proof}
If such a deduction existed it would apply to $\tilde\mu$ by Theorem~\ref{thm:mixture}(i)--(ii) and contradict (iii).
\end{proof}

\begin{remark}
Theorem~\ref{thm:mixture} sharpens the model of \cite[subsection ``Why these linear estimates leave the ordinary susceptibility open'']{dmn}, which separates linear-probe
concentration from quadratic susceptibility: here the obstruction law is itself a mixture of two actual Wilson laws,
so every \emph{averaged} full-law moment statement is inherited. The theorem does not say that the target is false
for $\mu_\beta$; it says which inputs a proof must use. It also explains why the residual analysis of
\cite[final subsection]{tra} could not close by adding further averaged moment estimates for the excess $E_f$: an averaged
moment bound is linear in the law.
\end{remark}

\section{The regime and the continuum trajectory}\label{sec:regime}

\begin{proposition}\label{prop:trajectory}
Let $(n_k,\beta_k)$ be a sequence with $n_k\to\infty$ and $b\log n_k\le\beta_k\le B\log n_k$ for some constants
$0<b\le B<\infty$ and all large $k$ (logarithmic running $\beta(a)\asymp\log(1/a)$ with $n=L/a$; \cite[subsection ``Mesoscopic scales remain mesoscopic'']{rsc}
and \cite{stp}). Then:
\begin{enumerate}
\item[(a)] $(n_k,\beta_k)\in\mathcal R$ for all large $k$;
\item[(b)] for every absolute constant $\beta_0$, $\beta_k\ge\beta_0$ for all large $k$; hence a hypothesis of the form
``$\beta\ge\beta_0$'' covers the trajectory;
\item[(c)] for every $\kappa>0$ and $C>0$, the hypothesis $\beta\ge Cn^{\kappa}$ fails along the trajectory for all large
$k$. In particular the flat-block threshold $\beta\ge C_0R^{14}\log(2+R)$ of \cite[block theorem of Section~2]{spw},
and the cube thresholds $\beta\ge CR^4$, $R\le(N\beta)^{1/48}$ of \cite[geometry-regime display and theorem on the complete averaged internal error]{cub}, permit only
$R\le C(\log n_k)^{1/14}$ and $R\le C(\log n_k)^{1/48}$ respectively, so that $aR\to0$: neither covers a fixed
physical block, as stated in \cite[subsection ``Mesoscopic scales remain mesoscopic'']{rsc}.
\end{enumerate}
No result proved in this note introduces any lower bound on $\beta$ beyond $\beta\ge1$. The missing inequality
(MI) below is stated with an absolute $\beta_0$, which is the only form compatible with (b)--(c).
\end{proposition}
\begin{proof}
(a) $c_*\log(2+N\beta_k)\le c_*\log(2+NB\log n_k)=o(n_k)$. (b) is immediate. (c) $Cn_k^\kappa>B\log n_k$ for large $k$.
The block-size statements follow by solving $C_0R^{14}\log(2+R)\le B\log n$ and $R^{48}\le NB\log n$ for $R$.
\end{proof}

\section{The isolated missing inequality and the deduction of the target}\label{sec:missing}

\subsection{Statement}

\begin{center}
\fbox{\parbox{0.92\textwidth}{
\textbf{(MI) Frustration structure-factor bound.} There exist constants $C_{\mathrm{MI}}<\infty$ and $\beta_0\ge1$,
depending only on $N$, $c_*$ and the fixed conventions, such that for all $(n,\beta)\in\mathcal R$ with
$\beta\ge\beta_0$, all $p\in\Lambda_n^*=\frac{2\pi}{n}\Lambda_n$ and all $v\in\mathbb C^4$ with $|v|=1$,
\[
\frac1{n^4}\,\E_{\mu_\beta}\Big|\sum_{x\in\Lambda_n}\sum_{\nu=1}^4v_\nu e^{ip\cdot x}\big(h_{(x,\nu)}-\E_{\mu_\beta}h_{(x,\nu)}\big)\Big|^2
\ \le\ C_{\mathrm{MI}},\qquad h_{(x,\nu)}=\frac{\beta}{3}\sum_{i<j}\big|F_{(x,\nu),i}-F_{(x,\nu),j}\big|^2 .
\]}}
\end{center}

(MI) is \emph{not} proved in this note. It is the exact content of the target, in the following sense.

\begin{theorem}[Deduction of the target from (MI)]\label{thm:deduction}
Assume (MI). Then for all $(n,\beta)\in\mathcal R$ with $\beta\ge\beta_0$ and all real marks $f$ (arbitrary sign and
support),
\[
\Var_{\mu_\beta}(r_f)\le\big(C_{\mathrm{MI}}+C_W\big)M_f,\qquad C_W=38C_2,
\]
and for the physical source $\Var_{\mu_\beta}\big(r_H(f)\big)\le c_\beta^{-2}a^{-4}(C_{\mathrm{MI}}+C_W)M_f$.
Conversely, if $\Var_{\mu_\beta}(r_f)\le CM_f$ for all real marks and all $(n,\beta)\in\mathcal R$ with
$\beta\ge\beta_0$, then (MI) holds with $C_{\mathrm{MI}}=C+C_W$.
\end{theorem}
\begin{proof}
(MI) says $\lam(\hat c^{\,h}(p))\le C_{\mathrm{MI}}$ for all $p$ (Theorem~\ref{thm:symbol}(a),(c) for $h$). By
Proposition~\ref{prop:ward-symbol}, $\lam(\hat c(p))\le C_{\mathrm{MI}}+C_W$, and Theorem~\ref{thm:symbol}(c) for $r$ gives
the first bound for all complex, hence all real, marks. The physical statement is the multiplication by
$c_\beta^{-2}a^{-4}$. The converse is the same argument in the other direction.
\end{proof}

\begin{corollary}[Closure of the other objectives under (MI)]\label{cor:closure}
Assume (MI). Then, in the common regime of \ref{I8} (periodic moment regime, embedded-cube hypotheses of
\cite[geometry-regime display]{cub} and its admitted range), each of the following holds with a constant depending only on
$C_{\mathrm{MI}},C_W$ and the constants of \ref{I8}:
\begin{enumerate}
\item[(1)] $\Var_{\mathcal B}\E(r_f\mid\mathcal B)\le CM_f$ for the constructed unrestricted cube observation
(uncontrolled-mean display of \cite{cub} and discarded-norm display of \cite{rsc});
\item[(2)] $\Var_{\tau,\mu}(K_\tau E_f)\le CM_f$ for the exact excess of \cite[residual display]{tra}, i.e.\ \cite[missing-bound display]{tra}, via the paid comparison of \cite{rsc};
\item[(3)] the source-space resolvent bound $A^*L_1^{-1}A\le CI$ of \cite[missing-bound display]{cur} for the specified first coordinate layer of
\cite[subsection ``An exact source-specific operator target'']{rsc}, via the resolvent identity of \cite{cur} and the complete innovation estimate \ref{I5};
\item[(4)] $\psi_f''(0)\le(C+C_6)M_f$ for all nonnegative marks, by \ref{I6}; conversely a curvature bound for
nonnegative marks gives the target for signed marks by Proposition~\ref{prop:signed} with constant $2(C+C_6)$.
\end{enumerate}
\end{corollary}
\begin{proof}
Theorem~\ref{thm:deduction} gives item~(1) of \ref{I8}; the remaining items are the equivalences of \ref{I8},
whose proofs (\cite[subsection ``Equivalence of the remaining quantitative objectives'']{rsc}) use only the exact variance decompositions, the paid comparison, the resolvent identity and \ref{I6}.
The signed-mark direction in (4) is Proposition~\ref{prop:signed}.
\end{proof}

\subsection{What is known about the symbol, and what is not}

\begin{proposition}\label{prop:known}
On $\mathcal R$:
\begin{enumerate}
\item[(a)] (Average.) $n^{-4}\sum_p\Tr\hat c^{\,h}(p)\le4C_2$ and $\#\{p:\lam(\hat c^{\,h}(p))>K\}\le4C_2n^4/K$
(Corollary~\ref{cor:mean}). The missing statement is the \emph{maximum} over $p$, not the mean.
\item[(b)] (Diagonal and near-diagonal.) $|c^h_{\nu\nu'}(z)|\le C_2$ for all $z$; and $c(z)=c^h(z)$ unless the
roots $(z,\nu),(0,\nu')$ share a plaquette, which requires $|z|_\infty\le1$ (Lemma~\ref{lem:locality}).
\item[(c)] (Tilt average.) For indicator marks, $\int_{-\delta/2}^{\delta/2}\Var_{\mu_\eta}(r_{\mathbf 1_A})d\eta\le C|A|$
(\cite[Part~VI, (T17)]{csj}); the value at $\eta=0$ is not controlled.
\item[(d)] (Gaussian reference.) In the Maxwell reference law of \cite[Part~III]{csj}, the kernel obeys
$|c^{(0)}(z)|\le C(1+|z|)^{-12}$ (\cite[Part~III, (2.7)]{csj}), hence its symbol is bounded by Proposition~\ref{prop:schur}; this is a
reference-law computation and not an actual-law comparison (\cite[subsection ``Relation to the unresolved compact susceptibility'']{gau}).
\item[(e)] (Flat block.) For the relative block law with identity tangential boundary and
$\beta\ge C_0R^{14}\log(2+R)$, \cite[block theorem of Section~2]{spw} gives $\Var(r_f)\le Cm_f$; by
Proposition~\ref{prop:trajectory}(c) this regime does not contain the periodic trajectory.
\item[(f)] (Obstruction.) By Theorem~\ref{thm:mixture}, (MI) is not a consequence of (a)--(c) together with the
component-wise Ward, contact, innovation and heat inputs.
\end{enumerate}
\end{proposition}
\begin{proof}
(a),(b) were proved above; (c)--(e) are citations with their stated scopes; (f) is Corollary~\ref{cor:obstruction}.
\end{proof}

\begin{remark}[What (MI) demands]
By Theorem~\ref{thm:symbol}(c), (MI) is the statement that the plane-wave frustration observables
$n^{-2}\sum_{x,\nu}v_\nu e^{ip\cdot x}(h_{(x,\nu)}-\E h_{(x,\nu)})$ have uniformly bounded second moment for every
momentum, including $p=0$ (the bulk susceptibility $n^{-4}\Var(\sum_{x}h_{(x,\nu)})$) and the soft momenta
$|p|\asymp n^{-1}$ singled out by the flat-toron calculation of \cite[proposition on a nonzero quadratic defect at a flat toron]{tra}. Under the exact score
identity, $\Cov(h_f,h_g)$ is a four-point force correlation whose Gaussian value is extensive; the nonlinear
control of its low-momentum part is exactly what \cite[Part~III, subsection ``What a nonlinear upgrade would have to \dots'']{csj} calls the missing
annealed low-spectral estimate. This note does not supply it.
\end{remark}

\section{The fixed-physical-volume form, the sourced pressure, and what the existence argument needs}\label{sec:fixedvolume}

The fixed-volume susceptibility criterion is needed only on the fixed reference torus $T^*$ of physical side $t^*$, i.e.\ for
$n=n^*(a)=t^*/a$ and $\beta=\beta(a)$ along the trajectory, as a corollary of the renormalization-group transport
of inserted observables (I-3(iv)). The volume-uniform formulation requires it for every $n\ge c_*\log(2+N\beta)$. This section
separates the two statements. Labels: (P) proved here, (L) literature, (C) conditional on a
displayed hypothesis, (O) open.

\subsection{The fixed-volume target from complex-source analyticity}

\begin{center}
\fbox{\parbox{0.92\textwidth}{
\textbf{Hypothesis $(\mathrm A_{T^*})$ (fixed-volume analyticity of the frustration-sourced pressure).}
There are $r_0>0$ and $C_A<\infty$ such that, for every lattice spacing $a$ on the admitted trajectory, with
$n=n^*(a)$ and $\beta=\beta(a)$, and for every \emph{real} mark $f$ on $\Lambda_{n}$ with $\|f\|_\infty\le1$, the entire
function $\varepsilon\mapsto Z_f(\varepsilon)=\E_{\mu_\beta}e^{\varepsilon h_f}$ has no zero in the closed disc
$\{|\varepsilon|\le r_0\}\subset\mathbb C$, and its logarithm (the branch with $\log Z_f(0)=0$) satisfies
$|\log Z_f(\varepsilon)|\le C_An^4$ there.}}
\end{center}

$(\mathrm A_{T^*})$ is the last clause of interface I-1 (``$p_{a,T^*}[J]$ analytic in $J$ on a fixed polydisc'')
specialized to the bare source $J=\varepsilon\sum_bf_bh_b$; it is also the ``zero-free complex source disc with
$|\log Z_H|\le Cn^4$'' named as a sufficient substitute in \cite[section ``A narrower route to the main susceptibility, and its current limit'']{rsm}.
It is a Lee--Yang-type statement and is \textbf{(O)}.

\begin{theorem}[(C) on $(\mathrm A_{T^*})$: the fixed-volume target]\label{thm:fixedvolume}
Assume $(\mathrm A_{T^*})$. Then for every $a$ on the trajectory, with $n=n^*(a)$, $\beta=\beta(a)$,
\[
\max_{p\in\Lambda_n^*}\lam\big(\hat c^{\,h}(p)\big)\le\frac{4C_A}{r_0^2},\qquad
\Var_{\mu_\beta}(r_f)\le\Big(\frac{4C_A}{r_0^2}+C_W\Big)M_f\quad\text{for every real mark }f,
\]
and $\Var(r_H(f))\le c_\beta^{-2}a^{-4}(4C_A/r_0^2+C_W)M_f$ for the physical source. No lower bound on $\beta$ beyond the
trajectory itself is used.
\end{theorem}
\begin{proof}
Fix a real mark with $\|f\|_\infty\le1$. On the finite lattice $h_f$ is bounded, so $Z_f$ is entire; under
$(\mathrm A_{T^*})$, $g=\log Z_f$ is analytic on $|\varepsilon|<r_0$ with $g(0)=0$, $g'(0)=\E h_f$ and
$g''(0)=\Var_{\mu_\beta}(h_f)$. Cauchy's estimate on the circle $|\varepsilon|=r_0$ gives
$\Var(h_f)=|g''(0)|\le 2r_0^{-2}\sup_{|\varepsilon|=r_0}|g(\varepsilon)|\le 2C_An^4/r_0^2$.
For a plane wave $f_{p,v}(x,\nu)=v_\nu e^{ip\cdot x}$ with $|v|=1$, both $\Re f_{p,v}$ and $\Im f_{p,v}$ are real marks with
sup norm at most one, and for complex marks $\Var(h_f)=\Var(h_{\Re f})+\Var(h_{\Im f})$ because the covariance form is
real symmetric. Hence $n^{-4}\Var(h_{f_{p,v}})\le4C_A/r_0^2$, which is the first bound by Theorem~\ref{thm:symbol}(c)
for $h$. Proposition~\ref{prop:ward-symbol} and Theorem~\ref{thm:symbol}(c) for $r$ give the second; the physical
statement is the multiplier.
\end{proof}

\begin{remark}[Interface change request on I-1]\label{rem:ICR}
Interface I-1 as written concerns smeared bare observables at the physical scale $\ell_0$ and chessboard-replicated
sources. Theorem~\ref{thm:fixedvolume} needs the analyticity clause for the \emph{bare lattice-scale} source
$\varepsilon\sum_bf_bh_b$ with arbitrary bounded real marks, including marks oscillating at the lattice scale, on the
polydisc $|\varepsilon|\le r_0$ with the bound $C_An^4$. This is the exact extra statement to be added to I-1
(\textbf{ICR-1}). It cannot be replaced by a bound for replicated or smooth marks alone, because
Theorem~\ref{thm:symbol}(c) requires every momentum $p$.
\end{remark}

\begin{remark}[Why reflection positivity does not supply $(\mathrm A_{T^*})$]
The chessboard estimate bounds $|Z_f(\varepsilon)|$ from above by replicated-source pressures, for complex
$\varepsilon$ as well (with conjugation in odd cells). An upper bound on $|Z_f|$ does not control $g''(0)$: Cauchy applied
to $Z_f$ instead of $\log Z_f$ gives only $|Z_f''(0)|\le2e^{C n^4}/r_0^2$. The content of $(\mathrm A_{T^*})$ is the
\emph{lower} bound $|Z_f(\varepsilon)|\ge e^{-C_An^4}$ on the complex disc, i.e.\ zero-freeness, and for real
$\varepsilon$ convexity of $\log Z_f$ with bounded values permits a narrow peak of $g''$ at zero (\cite[Part~VI, remark after (T17)]{csj}). This is the precise sense in which
``RP + local moments + real pressure stability'' do not suffice, and it is why $(\mathrm A_{T^*})$ is an
output of the renormalization-group transport (I-1) and not of the functional inequalities among the imported inputs.
\end{remark}

\subsection{The volume-uniform target is an infrared clustering statement that the existence argument does not need}

For $T=qT^*$ write $n=qn^*$ and let $\{c\}$ be the $q^4$ cells of side $t^*$ (lattice side $n^*$). For a fixed
orientation $\nu$ put $R_c=\sum_{x\in c}r_{(x,\nu)}$, the ordinary source summed over a cell.

\begin{theorem}[(P)]\label{thm:IR}
\begin{enumerate}
\item[(a)] (Cell structure factor.) If $\Sigma(qn^*,\beta)\le K$ then the $q^4\times q^4$ cell covariance matrix
$G_{cc'}=\Cov(R_c,R_{c'})$ satisfies $0\le G\le K\,n^{*4}\,I$ as a quadratic form on $\ell^2(\text{cells})$; in particular
$\Var(\sum_cg_cR_c)\le Kn^{*4}\sum_cg_c^2$ for all cell weights $g$, uniformly in $q$.
\item[(b)] (Bounded correlations do not imply (a).) Let $\mathcal L$ be any class of probability laws on the link
configurations of $T=qT^*$ that is closed under convex combinations and contains the Wilson laws $\mu_\beta$ and
$\mu_{\Lambda\beta}$ of Theorem~\ref{thm:mixture}; for example, the class of laws satisfying, for all $m$, bounds
$|\E\prod_{j=1}^mA_j|\le C^mm!\prod_j\|A_j\|$ for cell-supported insertions $A_j$ with constants independent of $q$.
Then $\mathcal L$ contains a law $\tilde\mu$ for which $\Var_{\tilde\mu}(\sum_cR_c)\ge\tfrac9{16}C_F^2N^2\,n^8$, so that
(a) fails for $\tilde\mu$ with any $K$ independent of $q$.
\item[(c)] (Not needed for reconstruction.) The Osterwalder--Schrader reconstruction in its corrected form requires
E0$'$--E3 and not the cluster property E4 (Osterwalder--Schrader II, CMP 42 (1975), Introduction Remark~6 and Section~IV;
Osterwalder--Schrader I, CMP 31 (1973), Remark~3, p.~89), and the Driessler--Summers decomposition (Helv.\ Phys.\ Acta 59
(1986), Theorem~2.3) produces components with unique vacuum without a rate of clustering. Hence no interface of the
the stated hypotheses requires $\Sigma(qn^*,\beta(a))$ to be bounded uniformly in $q$.
\end{enumerate}
\end{theorem}
\begin{proof}
(a) For cell weights $g$ take the mark $f(x,\nu)=g_c$ for $x\in c$ (one orientation); then $M_f=n^{*4}\sum_cg_c^2$ and
$\Var(r_f)=\sum_{c,c'}g_cg_{c'}G_{cc'}$. Apply the definition of $\Sigma$.
(b) Bounds of the displayed type are linear in the law and hold for $\tilde\mu=\frac12\mu_\beta+\frac12\mu_{\Lambda\beta}$
whenever they hold for both components, so $\tilde\mu\in\mathcal L$. Theorem~\ref{thm:mixture}(iii) with the mark
$f\equiv1$ on one orientation ($\|f\|_1=n^4$) gives $\Var_{\tilde\mu}(\sum_cR_c)=\Var_{\tilde\mu}(r_f)\ge\tfrac9{16}C_F^2N^2n^8$.
(c) is the cited literature.
\end{proof}

\begin{remark}
Theorem~\ref{thm:IR}(a) says that the volume-uniform target of \cite{cur} contains a uniform bound on the structure
factor of the \emph{ordinary source summed over cells of fixed physical size}, i.e.\ the absence of superextensive
fluctuations of block energies in arbitrarily large volumes: a weak infrared clustering property. (For Wilson laws
$\Sigma$ and $\Sigma^h$ differ by at most $C_W$, Proposition~\ref{prop:ward-symbol}; for the mixture law of (b) the
Ward reduction is not available, which is why (a) and (b) are stated for $r$.) No theorem in the literature gives this for
four-dimensional Yang--Mills at weak coupling, and by (b) it is not a consequence of uniformly bounded correlation
functions. By (c) it is also not required for Theorem YM-E of the stated hypotheses, which needs only the fixed-volume
statement, i.e.\ (MI) with $n=n^*(a)$. Accordingly, the fixed-volume target follows from
$(\mathrm A_{T^*})$, the strengthening ICR-1 of I-1, by Theorem~\ref{thm:fixedvolume}, and the
volume-uniform statement is a separate open problem.
\end{remark}

\subsection{Interface statement}

\emph{Assumed.} For Sections~\ref{sec:frustration}--\ref{sec:missing} and
Theorem~\ref{thm:IR}, only the imported inputs \ref{I1}--\ref{I8} with their stated locations. Hypothesis
$(\mathrm A_{T^*})$, equal to I-1's analyticity clause for bare lattice-scale frustration sources (ICR-1), for
Theorem~\ref{thm:fixedvolume} only.

\emph{Proved.} (P): Lemmas~\ref{lem:moments}--\ref{lem:locality}, Theorem~\ref{thm:symbol},
Proposition~\ref{prop:ward-symbol}, Corollaries~\ref{cor:planewave}--\ref{cor:mean},
Propositions~\ref{prop:schur}--\ref{prop:indicator}, Theorem~\ref{thm:mixture}, Corollary~\ref{cor:obstruction},
Proposition~\ref{prop:trajectory}, Theorems~\ref{thm:deduction} and~\ref{thm:IR}. (C): Theorem~\ref{thm:fixedvolume}
on $(\mathrm A_{T^*})$. (O): (MI), both in its fixed-volume form and in its volume-uniform form; $(\mathrm A_{T^*})$.

\emph{Required strengthening.} ICR-1 on I-1: analyticity and the bound $|\log Z_f(\varepsilon)|\le C_An^4$ on
$|\varepsilon|\le r_0$ for the bare source $\varepsilon\sum_bf_bh_b$, all real marks with $\|f\|_\infty\le1$, on $T^*$,
uniformly along the trajectory.

\section{Status}

\begin{itemize}
\item The target $\Var_{\mu_\beta}(r_f)\le CM_f$ is \textbf{not proved}. Its exact remaining content is (MI).
On the fixed reference torus it follows from hypothesis $(\mathrm A_{T^*})$ (Theorem~\ref{thm:fixedvolume}), which is
the analyticity clause of interface I-1 for bare frustration sources (ICR-1) and is open. Uniformly in the volume it
contains an infrared clustering statement (Theorem~\ref{thm:IR}) that the existence argument does not need.
\item Proved here, unconditionally and with constants depending only on $N$, $c_*$ and the fixed conventions:
Lemmas~\ref{lem:moments}--\ref{lem:locality}, Theorem~\ref{thm:symbol}, Proposition~\ref{prop:ward-symbol},
Corollaries~\ref{cor:planewave}--\ref{cor:mean}, Propositions~\ref{prop:schur}--\ref{prop:indicator},
Theorem~\ref{thm:mixture}, Corollary~\ref{cor:obstruction}, Proposition~\ref{prop:trajectory}, the
deduction Theorem~\ref{thm:deduction} with Corollary~\ref{cor:closure}, and Theorem~\ref{thm:IR}.
\item No additional lower bound on $\beta$ is introduced by the proved results. (MI) is formulated with an absolute
$\beta_0$; Proposition~\ref{prop:trajectory} shows this is the only admissible form along the logarithmic trajectory.
\item The flat-block restriction $\beta\ge CR^{14}\log(2+R)$, the separate cube restrictions, all raw normalizers
$Z_s,Z_j$, the raw-denominator responses, directional exterior curl, moving transports and shared-history means of
\cite{bhi,wwa,inn,whe,dmn,cur,tra,cub,gau,rsc} remain in force and untouched. Physical variance bounds carry $c_\beta^{-2}a^{-4}$. No physical-scale
induction by tiling is claimed, and no part of the Yang--Mills existence theorem beyond the stated reductions is
claimed.
\end{itemize}

\begin{thebibliography}{99}
\bibitem{cur} \emph{Conditional current estimates}, \texttt{lattice\_theory/current.tex} in this collection.
\bibitem{tra} \emph{Transport of the ordinary source}, \texttt{lattice\_theory/transport.tex} in this collection.
\bibitem{cub} \emph{Cube estimates for the ordinary law}, \texttt{lattice\_theory/cubes.tex} in this collection.
\bibitem{gau} \emph{Gaussian comparison of ordinary fluctuations}, \texttt{lattice\_theory/gaussian.tex} in this collection.
\bibitem{rsc} \emph{Conditional variance and the retained susceptibility},
\texttt{lattice\_theory/block\_and\_transfer/ordinary\_law/retained\_susceptibility\_criteria.tex} in this collection.
\bibitem{whe} \emph{Ward identities for the heat bath}, \texttt{lattice\_theory/ward\_heat.tex} in this collection.
\bibitem{dmn} \emph{Directional means and temporal sources}, \texttt{lattice\_theory/directional\_mean.tex} in this collection.
\bibitem{bhi} \emph{Boundary histories and conditional covariance}, \texttt{lattice\_theory/boundary\_history.tex} in this collection.
\bibitem{wwa} \emph{Ward identities for ordinary sources}, \texttt{lattice\_theory/ward\_ward.tex} in this collection.
\bibitem{inn} \emph{Innovation derivatives in boundary coordinates}, \texttt{lattice\_theory/innovation.tex} in this collection.
\bibitem{csj} \emph{Projected cross terms and compact retained-source jets},
\texttt{lattice\_theory/block\_and\_transfer/ordinary\_law/compact\_source\_jets.tex} in this collection.
\bibitem{spw} \emph{Spatial estimates for the original Wilson sources},
\texttt{lattice\_theory/block\_and\_transfer/conditional\_estimates/spatial\_estimates\_for\_the\_original\_wilson\_sources.tex} in this collection.
\bibitem{rsm} \emph{Rare selector means and inherited coefficients},
\texttt{lattice\_theory/block\_and\_transfer/conditional\_estimates/rare\_selector\_means\_and\_inherited\_coefficients.tex} in this collection.
\bibitem{sbr} \emph{Actual multi-plaquette Wilson tails and a sparse exceptional-slice record},
\texttt{lattice\_theory/block\_and\_transfer/wilson\_transfer/original\_sparse\_bath\_records.tex} in this collection.
\bibitem{stp} \emph{Original-law probability of an all-time staple tube},
\texttt{lattice\_theory/block\_and\_transfer/wilson\_transfer/original\_law\_small\_tube\_probability.tex} in this collection.
\end{thebibliography}

\end{document}
